\chapter*{常见缩写和术语先导}
\addcontentsline{toc}{chapter}{常见缩写和术语先导}

\begin{center}
\small
\begin{tabular}{@{}p{0.20\linewidth}p{0.32\linewidth}p{0.38\linewidth}@{}}
\toprule
缩写 & 全称 & 本卷中的含义 \\
\midrule
AI & Artificial Intelligence & 这里特指本地推理工程，不讨论泛泛应用口号。 \\
LLM & Large Language Model & decoder-only 文本生成模型是本卷主线。 \\
KV Cache & Key/Value Cache & attention 历史 K/V 状态，属于请求生命周期。 \\
GEMM & General Matrix Multiply & 矩阵乘，是预填充和训练常见核心算子。 \\
GEMV & General Matrix Vector Multiply & 矩阵向量乘，是单 token decode 的核心形态。 \\
SIMD & Single Instruction Multiple Data & 一条指令处理多个 lane 的 CPU 执行方式。 \\
AVX2 & Advanced Vector Extensions 2 & x86-64 上的 256-bit SIMD 指令集之一。 \\
FMA & Fused Multiply-Add & 乘加融合指令，常出现在浮点 kernel 中。 \\
SLO & Service Level Objective & 服务目标，例如 TTFT、TPOT、p95/p99。 \\
TTFT & Time To First Token & 从请求进入到首 token 输出的时间。 \\
TPOT & Time Per Output Token & decode 阶段每个输出 token 的平均或尾部时间。 \\
IR & Intermediate Representation & 编译器内部表示，用于 pass、lowering 和计划生成。 \\
RAG & Retrieval-Augmented Generation & 检索增强生成，需要把外部文档和 prompt 编排起来。 \\
MoE & Mixture of Experts & 多专家模型，会改变路由、负载和内存账本。 \\
LoRA & Low-Rank Adaptation & 低秩适配权重，常用于微调和多租户加载。 \\
\bottomrule
\end{tabular}
\end{center}

这些缩写不要死记。每个缩写在正文里都要落到一个代码对象或报告字段：KV cache 对应 trace 里的 slot 和 position，SLO 对应 serving probe 的 CSV，IR 对应 executable plan，SIMD/AVX2 对应 kernel 反汇编和 benchmark。
